% -*- coding: utf-8 -*-
\chapter{国内外研究现状与理论基础}
\label{chap:related_work}

可重构智能表面（RIS）作为突破 6G 高频通信覆盖瓶颈与能效约束的核心使能技术，其研究涵盖了电磁超材料物理建模、主被动协同波束赋形、高维级联信道获取以及实际硬件非理想损耗补偿等多个维度~\cite{wang2023road}。本章系统综述国内外在上述关键领域的研究演进与学术前沿，并深入剖析现有技术方案的局限性。

\section{RIS 电磁物理调控机理与信道建模研究}

传统无线通信通常假定无线信道是由散射体、反射体与移动障碍物构成的被动随机黑盒系统。RIS 的诞生使得收发机具备了主动调控空间电磁波传输路径的物理维度~\cite{wu2020towards}。从微观电磁学原理来看，RIS 表面由周期性排列的超材料反射单胞（Meta-atom）构成。每个单胞通常集成低功耗半导体器件（如 PIN 二极管或变容二极管）。通过控制加载到二极管两端的偏置直流电压，能够改变微元的表面等效电感与电容参数，进而产生不同的谐振响应，在宏观上表现为对入射电磁波幅度和相位的灵活控制。

在早期的理论分析中，文献通常将 RIS 简化为理想的连续纯相移对角矩阵模型，即假定各反射单元的幅度反射系数 $\beta_n = 1$，且相移角 $\theta_n$ 可在区间 $[0, 2\pi)$ 内连续任意取值。然而，在实际硬件工程实现中，微元的相位调整往往伴随着不可避免的电磁寄生电阻损耗。文献~\cite{shen2022modeling}基于多端口散射参数网络分析方法，建立了物理上自洽的非理想 RIS 硬件耦合模型，明确指出了反射相位与幅度之间的内在物理耦合关系（即相移从 $0$ 变至 $\pi$ 时伴随幅度反射跌落）。同时，为了降低 FPGA 控制器的硬件复杂度与控制总线功耗，实际商用 RIS 多采用有限精度的离散相移量化（如 1-bit 或 2-bit）。文献~\cite{wu2020towards}在 MISO 架构下的系统分析表明，当相移量化阶数达到 3-bit（8 个离散相位状态）时，离散控制即可逼近连续理想相移 99\% 以上的等效信噪比增益。

\section{主被动协同波束赋形优化算法演进}

主被动协同波束赋形是发挥 RIS 空间调控潜能的核心算法支撑。在典型多天线基站（配备 $M$ 根天线）与多用户通信系统中，基站数字预编码向量 $\bm{W}$ 与 RIS 反射相移向量 $\bm{v}$ 的优化目标通常为最大化系统总可达和速率、最小化基站发射功率或最大化全网能量效率。

\subsection{经典交替优化与半正定松弛算法}
文献~\cite{wu2019intelligent}首次提出了基于交替优化（Alternating Optimization, AO）的主被动波束赋形框架。该方法将高维非凸耦合问题分解为两个交替迭代的子问题：在固定 RIS 相移矩阵 $\bm{\Theta}$ 时，基站发射波束赋形转化为标准的最大比合并（MRC）或迫零（ZF）凸优化问题；在固定基站波束时，针对包含单位模约束 $|e^{j\theta_n}| = 1$ 的非凸 RIS 相移子问题，采用半正定松弛（Semidefinite Relaxation, SDR）技术将秩一非凸约束松弛为半正定凸优化，并借助高斯随机化还原可行解。研究证明，在理想视距信道下，RIS 能够提供随反射单元数平方增长的功率缩放定律（$\mathcal{O}(N^2)$ 功率增益）。然而，SDR 算法升维带来的高额计算复杂度使得其难以拓展至超大规模阵列。

\subsection{双时间尺度与非理想 CSI 鲁棒优化}
针对实际系统中信道测量开销过高的问题，文献~\cite{zhi2023timescale}提出了双时间尺度（Two-Timescale）优化机制。该机制利用信道的大尺度统计特性优化变动较慢的 RIS 反射相位，而仅在小尺度瞬时相干时间内快速调节基站基带预编码，显著降低了信道反馈频次。文献进一步证明，在统计信道信息（CSI）不完备条件下，双时间尺度设计仍能逼近完全瞬时 CSI 方案 90\% 以上的遍历容量。

\subsection{深度强化学习驱动的自适应波束调控}
面对复杂多变的高速移动与动态干扰场景，传统数学规划算法依赖精确先验信道模型且单次迭代耗时较高。文献~\cite{yang2021deep}引入了深度确定性策略梯度（DDPG）无模型强化学习算法。智能体将环境瞬时特征映射为状态向量，通过连续动作网络直接生成微秒级连续相移输出，实现了在动态快衰落与潜在物理层窃听环境下的自适应快速决策。然而，强化学习方法属于黑盒优化，训练收敛性依赖大量样本探索，在通信系统的严苛确定性保障方面仍存在理论短板。

\section{无源 RIS 级联信道估计关键技术}

由于无源 RIS 不配备有源射频接收链，发射信号经历“基站-RIS-用户”两次自由空间衰落后的级联信道获取成为业界公认的“瓶颈中的瓶颈”。

早期文献采用基于 ON/OFF 开关协议的穷举信道估计法，即依次开启单个 RIS 反射微元而关闭其余所有微元，从而独立测量各级联子信道。这种方案实现简单但信噪比损失严重（关闭绝大多数微元损失了反射阵列增益），且导频训练长度与单元数 $N$ 严格成正比。文献~\cite{he2021channel}针对毫米波 RIS MIMO 系统信道的空间角域稀疏性，提出了基于原子范数最小化（Atomic Norm Minimization, ANM）的无网格信道估计算法，利用压缩感知技术突破了传统离散网格的量化误差限制。文献~\cite{liu2020matrix}进一步针对多用户下行信道，提出了基于矩阵校准（Matrix Calibration）的双线性最小二乘估计框架，利用不同用户与基站间共享的基站-表面信道矩阵内在低秩结构，大幅压缩了系统导频开销。

\section{代表性文献综合对比分析}

为了清晰展示本论文课题在学术脉络中的定位，基于本地真实文献库提炼的代表性核心文献多维对比详见表~\ref{tab:ris_literature_comparison}。

\begin{table}[htbp]
  \centering
  \small
  \caption{面向 6G 的 RIS 辅助通信波束赋形与信道估计代表性文献对比表}
  \label{tab:ris_literature_comparison}
  \begin{tabular*}{\textwidth}{@{\extracolsep{\fill}}clllllc@{}}
    \toprule
    \textbf{序号} & \textbf{代表文献} & \textbf{年份/期刊} & \textbf{核心技术范式} & \textbf{核心增益与评估指标} & \textbf{与本课题关系} & \textbf{重点精读} \\
    \midrule
    1 & Wang et al.~\cite{wang2023road}     & 2023 IEEE COMST & 6G 全景愿景与使能技术 & 频谱能效提升 10$\sim$100 倍 & 绪论立项依据 & \textbf{是} \\
    2 & Wu \& Zhang~\cite{wu2019intelligent}& 2019 IEEE TWC   & 主被动波束交替优化(AO) & $\mathcal{O}(N^2)$ 渐近功率增益 & 理论核心基石 & \textbf{是} \\
    3 & Liu et al.~\cite{liu2020matrix}     & 2020 IEEE JSAC  & 级联信道矩阵校准估计   & 导频开销压缩至最小正交维 & 信道估计支撑 & \textbf{是} \\
    4 & Wu \& Zhang~\cite{wu2020towards}    & 2020 IEEE MCOM  & 电磁调控与有限精度量化 & 3-bit 恢复 99\% 连续增益   & 硬件模型参考 & \textbf{是} \\
    5 & Yang et al.~\cite{yang2021deep}     & 2021 IEEE TWC   & DDPG 深度强化学习波束 & 安全速率提升 35\%          & 动态算法对比 & 否 \\
    \bottomrule
  \end{tabular*}
  \vspace{0.15cm}\par
  \noindent\footnotesize{
    \textbf{注}：本表全部数据与结论严格锚定于本地物理文献库 \texttt{ris\_thesis\_workflow/literature/} 中对应的真实 OA 全文 PDF 原文。
  }
\end{table}

\section{本章小结}

本章系统梳理了可重构智能表面（RIS）在电磁建模、主被动联合波束赋形与高维级联信道估计三大方向的国内外前沿研究进展：
\begin{enumerate}
  \item 分析了从理想连续相移到具有幅相耦合及离散量化特性的微元物理模型演进；
  \item 总结了以交替优化、半正定松弛及强化学习为代表的波束优化算法，指出了高维相移空间非凸性导致的求解复杂度困境；
  \item 剖析了无源 RIS 架构下由于缺乏射频链带来的级联信道估计超长导频开销难题。
\end{enumerate}
上述文献的深入分析表明，如何设计低复杂度、高收敛速率的主被动波束协同解耦算法，以及如何利用信道代数结构实现低开销正交导频信道估计，是当前 6G 物理层研究的迫切需求，直接构成了第三章算法设计的核心出发点。
